\documentclass[11pt]{article}
\usepackage[margin=25mm]{geometry}
\usepackage{fontspec,unicode-math}
\setmainfont{Latin Modern Roman}
\setmathfont{Latin Modern Math}
\usepackage{amsmath,amssymb,mathtools}
\usepackage{xurl,hyperref}
\hypersetup{colorlinks=true,urlcolor=blue,linkcolor=blue}
\setcounter{secnumdepth}{0}
\setlength{\parindent}{0pt}
\setlength{\parskip}{5pt}
\setlength{\emergencystretch}{3em}
\newcommand{\sref}[2]{\hyperlink{source:#1:#2}{[S#2]}}
\newcommand{\cataloguescope}[1]{\paragraph{Source question.}#1}
\newcommand{\floor}[1]{\left\lfloor#1\right\rfloor}
\begin{document}
\section{Randomized Communication Complexity of Distributed Triangle Detection}\label{statement}

\subsection{Definitions and mathematical statement}

\paragraph{Shared notation.}Standard number systems are denoted by $\mathbb N,\mathbb Z,\mathbb Q,\mathbb R,\mathbb C$; index ranges, computational models and specialized notation follow the cited problem source and the definitions below.


A triangle is three distinct vertices with all three joining edges. In the described vertex-partition model each participant knows incident edges, so an edge has two informed endpoint participants.

Communication complexity measures the number of transmitted bits needed for parties holding separate inputs to compute a joint function. Randomized protocols allow random choices, with their error probability specified by the model.

In the LOCAL model vertices communicate synchronously along graph edges; messages and local computation are unrestricted and complexity is the number of rounds. After $t$ rounds a vertex can depend only on its radius-$t$ neighborhood and any supplied identifiers or randomness.
\paragraph{Statement recorded in the supplied source.}
What is the randomized communication complexity of detecting a triangle when each graph vertex is assigned to a player who receives all incident edges, so each edge is known to its two endpoint players? \sref{statement}{2}
\paragraph{Scope qualification.}
The communication topology, public/shared randomness, and success probability are not specified.
\paragraph{Source status.}
The supplied catalogue labels this record partially solved. Its source-status metadata was checked against the supplied snapshot on 2026-10-11. This import does not independently establish current open status or mathematical correctness. The catalogue literature assessment is dated 2026-08-21.
\subsection{Short English statement}
Determine how much randomized communication is needed to detect a triangle when each participant sees one vertex' s edges.
\subsection{Sources}
\begin{itemize}
\item[\hypertarget{source:statement:1}{S1}] Ulam AI and the UnsolvedMath Contributors, UnsolvedMath: A Curated Collection of Open Mathematics Problems, catalogue entry. Catalogue curation and metadata: CC BY 4.0. The supplied source question is retained; definitions and the English summary are editorial additions.
\url{https://huggingface.co/datasets/ulamai/UnsolvedMath}.
\item[\hypertarget{source:statement:2}{S2}] Complexity Theory (2016), Oberwolfach Reports, DOI 10.4171/owr/2015/54, pp. 3049--3099
\url{https://doi.org/10.4171/owr/2015/54}.
\item[\hypertarget{source:statement:3}{S3}] Distributed triangle detection lower bounds
\url{https://arxiv.org/abs/2504.01802}.
\item[\hypertarget{source:statement:4}{S4}] FOCS 2025 version
\url{https://doi.org/10.1109/FOCS63196.2025.00040}.
\item[\hypertarget{source:statement:5}{S5}] Distributed triangle detection algorithms
\url{https://arxiv.org/abs/1807.06624}.
\textit{Underlying Oberwolfach Reports material retains its source-specific licence, including CC BY-SA 4.0 where applicable; see THIRD\_PARTY\_NOTICES.md.}
\end{itemize}
\paragraph{Preparation and review.}
Original question attributed to the cited source and UnsolvedMath contributors. Codex (GPT-6; exact serving revision not exposed) assisted with definitions, English summaries, source inspection and formatting. No new human mathematical authorship or independent proof verification is asserted. Catalogue identity checked against the complete supplied registry and existing repository statements; known flagged repetitions and obvious identical conjectures excluded. Full mathematical-equivalence/literature review has not been performed. Statement preparation and final review followed contribution prompts C and D; automated checks verify structure and identity, not open status.
\label{end:statement}
\end{document}
